% BEGIN RCB:R68_CONTINUATION
The complete singular-parameter connection acts on this same
rectangular image. Write \(f_i(s_i)=\Phi_h(s_i)-t_is_i\) and retain
the literal monic divisions
\[
 f_i(s_i)s_i^b=(h(s_i)-t_i)q_{i,b}+r_{i,b},\qquad
 (C_h(t_i))_b=r_{i,b},\qquad (B_h)_b=q_{i,b}' .
\]
The differential remainder is \(r_{i,b}-uq_{i,b}'\), with
\(C_h(t_i)=f_i(A_h(t_i))\). The original one-factor matrix is
\(\Omega_{u,h}(u,t_i)=-C_h(t_i)/u^2+B_h/u\), and
\(\partial_u\Pi_h=\Pi_h\Omega_{u,h}\).
The complete division and contour proof is in the boundary
connection chapter and (XD.u1)--(XD.u6), applied to the displayed
original \(h\)-variable. On the ordered product complex the
degree-\(j\) cochain operator is exactly
\[
 \nabla_u^{(j)}=\partial_u-\frac{\sum_i f_i}{u^2}+\frac{k-j}{u}.
\]
Its cochain identity follows one slot at a time from the
degree-zero/degree-one identity: increasing exterior degree
removes exactly one \(1/u\) term. Every exterior sign and the
original ordered form are retained.

Set \(\Pi_{\rm prod}=\bigotimes_i\Pi_h(u,t_i)\); let
\(\Omega_{u,\rm prod}\) be the sum of the \(k\) matrices
\(\Omega_{u,h}(u,t_i)\) acting in their respective tensor slots.
The product rule gives the exact original map
\[
 \partial_u\mathcal P
       =\Pi_{\rm prod}\Omega_{u,\rm prod}\eta_{h,k}.
\]
For the full original positive tensor coefficient Gram \(G\)
at the specified finite source cutoff, retain the explicit left
inverse
\[
 \ell_\eta=(\eta_{h,k}^*G\eta_{h,k})^{-1}\eta_{h,k}^*G,\qquad
 \ell_\eta\eta_{h,k}=I.
\]
The inverse exists because \(\eta_{h,k}\) is injective and \(G\)
is positive on the original full tensor source. The complete
parameter defect is therefore
\[
 \partial_u\mathcal P-\mathcal P
              (\ell_\eta\Omega_{u,\rm prod}\eta_{h,k})
   =\Pi_{\rm prod}(I-\eta_{h,k}\ell_\eta)
                  \Omega_{u,\rm prod}\eta_{h,k}.
\]
Substitution and factoring prove this equality, retaining the
actual transverse component and full unit in \(\eta_{h,k}\).
The observed form remains
\(\eta_{h,k}^*\Pi_{\rm prod}^*\Pi_{\rm prod}\eta_{h,k}\),
with its original source comparison Gram. The cyclic
\(\chi_{h,k}\)-potential has its own connection (XD.u1)--(XD.u7);
this injection, left inverse and transverse map give the precise
comparison with the product potential.

On the complete primary slice \(h=(s-\rho)^m\), \(t_i=0\), the
exact product period matrix, source-unit character projectors,
terminal arithmetic eigenclass and finite-field boundary are
proved in (BTP.1)--(BTP.11) and (BTF.1)--(BTF.11). The complex
invariant dimension is \((m^k+m(-1)^k)/(m+1)\); the actual
finite-field inertia dimension is that number when \(kc=0\)
and zero otherwise, with \(c=-(-\rho)^{m+1}/(m+1)\).
The original unit-weighted coefficient map and all its factorial
losses are (PC30a)--(PC30g) and (BPC.1)--(BPC.14), on precisely
their stated coefficient-map domains. These full proofs retain
the one-factor vanishing as its \(k=1\) case and carry each
surviving tensor character to its actual boundary stalk.
% END RCB:R68_CONTINUATION

% BEGIN RCB:R69_CONTINUATION
The same finite transport has its complete original
\(u\)-derivative. Keep \(\Omega_{u,h}=-C_h(t)/u^2+B_h/u\),
\(\Omega_{t,h}=-A_h(t)/u\), and \(t_b=t+ub\).
Differentiating the exact matrix formula at fixed \(b\) gives
\[
 \begin{aligned}
 \partial_u C_b&=-\Omega_{u,h}(u,t)C_b+
 C_b\bigl(\Omega_{u,h}(u,t_b)+b\Omega_{t,h}(u,t_b)\bigr),\\
 \partial_t C_b&=-\Omega_{t,h}(u,t)C_b+
                          C_b\Omega_{t,h}(u,t_b).
 \end{aligned}
\]
Indeed \(\partial_a\Pi^{-1}=-\Omega_{a,h}\Pi^{-1}\);
the target derivative is
\(\Pi(u,t_b)(\Omega_{u,h}+b\Omega_{t,h})(u,t_b)\).
This proves the pullback connection under the actual base map
\((u,t)\mapsto(u,t+ub)\), including \(b\partial_t\).
It is the polynomial restriction of the fully proved (APU23)
transport. These equations hold on the original nonzero-\(u\)
contour domain; the same entire finite-matrix ODE supplies the
regular transport through \(u=0\). The reflection and
inverse-transpose dual retain the stated parameter pullbacks
and original moment-line factor; their derivatives follow by
these same inverse and chain rules on the unchanged maps.
% END RCB:R69_CONTINUATION

% BEGIN RCB:R70_CONTINUATION
The formal gauge carries the complete regular form of the
singular-parameter operator. Put \(f=\Phi_h-ts\). On the original
two formal complex terms define
\[
 H_j=u^2\partial_u-f+(1-j)u,\quad j=0,1,\qquad
 H_1D=DH_0,\quad D=u\partial_s+h-t.
\]
The commutator of \(u^2\partial_u\) with \(D\) is
\(u^2\partial_s\), and that of \(-f\) with \(D\) is
\(u(h-t)\). Their sum is \(uD\), canceled by the extra \(u\)
on the degree-zero term. This proves the cochain identity.
These operators act coefficientwise on the stated completed
rings. After the specified localization in \(u\), they are
\(u^2\nabla_u^{(j)}\), with the same degree terms as
(APU1)--(APU4). Since \(\nu\) is independent of \(u,t\),
\[
 \nu H_j\nu^{-1}=H_j,\qquad
 H_1(\nu D\nu^{-1})=(\nu D\nu^{-1})H_0.
\]
These are operators on the transported gauge terms;
multiplication by \(f\) commutes with \(\nu\).
The polynomial-to-localized inclusion commutes with \(H_j\).
Its cohomology isomorphism, proved by the full contraction \(J\),
consequently has inverse commuting with the induced \(H\):
multiply the intertwining identity by that inverse on both
sides. This proves compatibility without assigning strict
horizontality to an uncomputed chain retraction. The exact
homotopy and full gauge endpoint remain (UG8)--(UG19).

The existing \(L=u\partial_t-s\) satisfies
\([H_j,L]=uL\), because its two nonzero terms are
\(u^2\partial_t\) and \(-us\). Including the derivative of
\(1/u\) therefore gives
\([H_j/u^2,L/u]=0\) on the localized complex.
At the original formal fibre \(H_1\) is multiplication by
\(-[\Phi_h]\), commuting with the entire \([\nu]=\upsilon_h\).
On the ordered \(k\)-fold complex the same calculation gives
\(u^2\partial_u-\sum_i(\Phi_h-t_is_i)+(k-j)u\);
the original signed tensor gauge and its source endpoint retain
their full degree terms and Taylor units.
% END RCB:R70_CONTINUATION

% BEGIN RCB:R71_CONTINUATION
The full connection of this analytic extension now includes its
\(u\) direction, as proved in (APU1)--(APU23). Write
\(f(s,t)=\Phi_h(s)-ts\), and let \(Q\) have columns
\[
 Q_p(s,t)=\frac{f(s,t)-f(p,t)}{s-p}
              -\frac{h(s)-t}{d+1},\qquad p\in\mathcal P_\nu.
\]
The first polynomial has degree \(d\) and leading coefficient
\(1/(d+1)\). Its division by \(h-t\) has that constant quotient;
the derivative correction is zero and the displayed remainder
has degree less than \(d\). In the original polynomial and
simple-pole columns the exact matrix is
\[
 \Omega_{u,\rm ext}=
 \begin{pmatrix}
 -C_h(t)/u^2+B_h/u&-Q/u^2\\
 0&-\operatorname{diag}(f(p,t))/u^2
 \end{pmatrix},\qquad
 \partial_u\Pi_{\rm ext}=\Pi_{\rm ext}\Omega_{u,\rm ext}.
\]
Apply \(\nabla_u=\partial_u-f/u^2\) to \(1/(s-p)\) and
separate \(f(p)/(s-p)\) from its complete polynomial divided
difference to prove this formula. The original \(t\)-matrix
remains
\(-u^{-1}\bigl(\begin{smallmatrix}A_h(t)&B_0\\0&P\end{smallmatrix}\bigr)\).
The full block flatness, including its off-diagonal block, is
proved in (APU8)--(APU12) using these literal remainders.

Put \(W_{\rm res}=\operatorname{diag}(e^{f(p,t)/u})\).
The weighted-residue row matrix is \(\mathcal R=[0,W_{\rm res}]\);
the positively oriented loop row matrix is \((2\pi i)\mathcal R\).
Differentiation gives
\[
 \partial_u W_{\rm res}
    =-\operatorname{diag}(f(p,t))W_{\rm res}/u^2,\qquad
 \partial_t W_{\rm res}=-PW_{\rm res}/u.
\]
Thus \(\partial_a\mathcal R=\mathcal R\Omega_{a,\rm ext}\)
for \(a=u,t\); both residue and loop maps are horizontal, with
their original \(2\pi i\) factors retained. The entire transport
has the two derivative identities in R69 with the full extension
matrices. Equation (APU23) proves the resulting shifted endpoint
and residue diagrams, including
\(b\Omega_{t,\rm ext}(u,t+ub)\).

Multiplication by the complete \(\nu\) commutes with the
parameter operators and transports the entire frame; its weight
is still \(w/\nu\). Its inverse-transpose dual and unchanged
endpoint \(J_0\) are proved in (APU18)--(APU22).
The original theta-source observation is exactly
\(\sigma_h[\nu]J_0\), retaining the entire unit.
When \(\nu\) is constant, \(r=0\) and every pole block, row
and additional loop is empty; the same equations reduce to the
full original polynomial connection and its unit gauge.
These maps retain the analytic family, the specified formal
endpoint and the formal coefficient operations of R70.
% END RCB:R71_CONTINUATION
